\scanpage{20}
- \mspage{16} -
\uline{Preuve} : (I) D'après (6.3), il n'y a en effet pas plus de deux automorphismes.

(II) et (III)\quad Lorsque les courbes sont mises sous la forme (2.1), il résulte de (2.5) que les automorphismes sont décrits parleur action sur $\pi$ et que $\lambda\in G_m$ défini un automorphisme si et seulement si
\[
c_4=\lambda^4c_4\qquad\text{et}\qquad c_6=\lambda^6c_6\qquad,
\]
d'où les assertions.

(IV)\quad Si la courbe est mise sous forme (6.1), les automorphismes sont donnés par (1.2) : $x'\leadsto x$, pour $s=t=0$ et
\[
\left\{\begin{aligned}
u^4&=1\\
r^3+ra_4+a_6(1-u^2)&=0
\end{aligned}\right.
\]
et la loi de composition par
\[
(u',r')\quad(u,r)=(u\,u',\ u'^2r+r')
\]
Lorsque la courbe est mise sous forme (6.7), ces équations se réduisent à $u\in\mu_4$ et $r\in\mathrm F_3$.

(V)\quad Si la courbe est mise sous forme (6.3), les automorphismes sont donnés par (1.2) : $x'\leadsto x$, pour $r=s^2$ et
\[
\left\{\begin{aligned}
u^3&=1\\
s^4+sa_3+a_4(1-u)&=0\\
t^2+ta_3+s^6+s^2a_4&=0
\end{aligned}\right.
\]
\margin{19}
